\documentclass[10pt,a4paper]{article}
\usepackage{amsmath,amssymb,amsthm,graphicx}
\usepackage{hyperref}
\usepackage{latexsym}
\usepackage{amssymb}
\usepackage{amsthm}
\usepackage{ulem}
\usepackage{enumitem}
\setlist{itemsep=0.4pt,topsep=0pt}
\parindent 0pt
\parindent 0pt
\parskip 12pt
\newcommand{\HCF}{\mathrm{hcf}}
\newcommand{\ZZ}{\mathbb{Z}}
\newcommand{\rk}{\mathrm{rk}}
\DeclareMathOperator{\len}{length}
\DeclareMathOperator{\bl}{b}
\DeclareMathOperator{\cl}{cl}
\DeclareMathOperator{\di}{d}
\DeclareMathOperator{\Cl}{Cl}
\DeclareMathOperator{\acl}{acl}
\DeclareMathOperator{\tp}{tp}
\DeclareMathOperator{\Th}{Th}
\newcommand{\Cor}{\mathcal{C}_\mathrm{or}}
\newcommand{\Cun}{\mathcal{C}_\mathrm{un}}
\newcommand{\mons}{\mathbb{M}}
\newcommand{\strong}{\mathrm{strong}}
\newcommand{\axioms}{\Delta_{\mathrm{FSTS}}}
\newtheoremstyle{note}% hnamei
{8pt}% hSpace abovei
{14pt}% hSpace belowi
{\normalfont}% hBody fonti
{}% hIndent amounti1
{\bfseries}% hTheorem head fonti
{ }% hPunctuation after theorem headi
{.5em}% hSpace after theorem headi2
{}% hTheorem head spec (can be left empty, meaning ‘normal’)i
\theoremstyle{note}
\newtheorem{thm}{Theorem}[section]
%\newtheorem{theorem}{Theorem}
\newtheorem{lemma}[thm]{Lemma} 
\newtheorem{df}[thm]{Definition}
\newtheorem{rmk}[thm]{Remark} 
\newtheorem{prop}[thm]{Proposition}
\newtheorem{fact}[thm]{Fact}
\newtheorem{cor}[thm]{Corollary}
\newtheorem{conj}[thm]{Conjecture}
\newtheorem{ex}[thm]{Example}
\newtheorem{qu}[thm]{Question}
\newtheorem*{notation}{Notation} 
\title{An infinite system has no maximum in an HF-ordering and every element lies below a product}
\author{Silvia Barbina, Enrique Casanovas}
\date{arxiv:2411.13723v3; Creative Commons Attribution 4.0 International}
\begin{document}
\maketitle
\noindent\textit{Editorial note.} The counted statement is the author's environment at source0.tex lines 829-834 and its proof environment opens at line 835; the pairing is by adjacency and the retained proof range is the whole environment.\par\medskip
\noindent\textit{Editorial note.} The preamble is the author's own package lines, macro block and theorem-environment declarations (source lines 4-65); the author's optioned amsmath line and the fleqn-dependent layout line are not retained.\par\medskip
\noindent\textit{Editorial note.} The source carries no label-number dump and no counter values are imposed, so the printed numbering is the wrapper's default; identity is the author label, or the pinned source position where the author gave none.\par\medskip
\noindent\textit{Editorial note.} Metadata provenance: the source record of the pinned copy and the licence evidence file are both present; the licence is CC BY 4.0 with status supported, and the card's source\_version matches the value recorded in that evidence file.\par\medskip
\noindent\textit{Editorial scope.} One original result of the pinned author manuscript, counted once: no maximum in an HF-ordering of an infinite system, and products above any prescribed element (author label tree2).\par\medskip
\noindent\textit{Allowed standard prerequisites (editorial).}\begin{itemize}
\item HF-ordering of an infinite system. A linear ordering in which each element has at most one pair of smaller elements forming a block with it.. It is the ordering under consideration in both items.
\item Block. A block is a triple of elements closed under the partial operation of the Steiner triple system.. It is what the maximality argument counts.
\item Infinite Steiner triple system. The system has more than three elements, so each element lies in several blocks.. It is what makes the first item work.
\end{itemize}\medskip
\begin{lemma}\label{tree2}
	\begin{enumerate}
		\item An HF-ordering of an infinite STS does not have a maximum.
		\item If $<$ is an HF-ordering of an infinite Steiner triple system $M$, then for every $a\in M$, for every  $b \geq a$, there are $b_1,b_2>b$ such that  $a=b_1\cdot b_2$.
	\end{enumerate}
\end{lemma}

\begin{proof}
	1. Assume $a\in M$ is a maximum. Since $|M|>3$, the element $a$ belongs to two different blocks, $\{a,a_1,a\cdot a_1\}$ and $\{a,a_2,a\cdot a_2\}$.  By maximality, $a> a_1,a_1\cdot a, a_2, a_2\cdot a$, in contradiction with $<$ being an HF-ordering.
	
	2. Assume not, and take $b \geq a$ such that  if $a=a_1\cdot a_2$, then $a_1\leq b$ or $a_2\leq b$. Since $M$ is infinite, by part 1 we can choose $b_1>b_2>b$ such that  $b_1\cdot b_2<b_1$. Note that $\{b_1,a,a\cdot b_1\}\neq \{b_1,b_2, b_1\cdot b_2 \}$, since otherwise $a=b_1\cdot b_2$ and  $b_1,b_2>b$. If $b_ 1> a\cdot b_1$, then  the previous blocks contradict that $<$ is an HF-ordering. Hence, $b_1<a\cdot b_1$. In this case, $a= b_1\cdot (a\cdot b_1)$  with $b_1,a\cdot b_1>b$, a contradiction with the choice of $b$.
\end{proof}

\end{document}
